\section{The Resumed Sundays after the Epiphany}
\label{sec:resumed}

\subsection{The rule at its own locus}

The mechanism has three loci, and they are not redundant: one states the arrangement, one
states the chants, and one states the trigger. Take them in that order.

\textbf{The arrangement} is \rn{18} of the \latin{Rubricae generales}:

\begin{studyframe}
\latin{18. Dominicae post Epiphaniam quae, superveniente Septuagesima, impediuntur,
transferuntur post dominicam XXIII post Pentecosten, hoc ordine:}
\begin{enumerate}[label=\textit{\alph*})~, leftmargin=1.9em, itemsep=0.15em, topsep=0.25em]
\item \latin{si dominicae post Pentecosten fuerint 25, dominica XXIV erit quae inscribitur
dominica VI post Epiphaniam;}
\item \latin{si dominicae fuerint 26, dominica XXIV erit quae inscribitur V post Epiphaniam;
et XXV, quae inscribitur VI;}
\item \latin{si dominicae fuerint 27, dominica XXIV erit quae inscribitur IV post Epiphaniam;
XXV, quae inscribitur V; et XXVI, quae inscribitur VI;}
\item \latin{si dominicae fuerint 28, dominica XXIV erit quae inscribitur III post Epiphaniam;
XXV, quae inscribitur IV; XXVI, quae inscribitur V; et XXVII, quae inscribitur VI.}
\end{enumerate}
\latin{Ultimo tamen loco semper ponitur ea quae in ordine est XXIV post Pentecosten, omissis, si
opus sit, ceteris, quae aliquando locum habere non possunt.}
\medskip

\textit{Working gloss.} The Sundays after the Epiphany that are impeded by the coming of
Septuagesima are transferred after the Twenty-third Sunday after Pentecost, in this order:
\textit{a)} if there are 25 Sundays after Pentecost, the Twenty-fourth will be the one inscribed
as the Sixth after the Epiphany; \textit{b)} if there are 26, the Twenty-fourth will be the one
inscribed as the Fifth after the Epiphany, and the Twenty-fifth the one inscribed as the Sixth;
\textit{c)} if there are 27, the Twenty-fourth will be the Fourth after the Epiphany, the
Twenty-fifth the Fifth, and the Twenty-sixth the Sixth; \textit{d)} if there are 28, the
Twenty-fourth will be the Third after the Epiphany, the Twenty-fifth the Fourth, the
Twenty-sixth the Fifth, and the Twenty-seventh the Sixth. In the last place, however, is always
put the one that is in order the Twenty-fourth after Pentecost, the rest being omitted, if need
be, which sometimes cannot have a place.
\end{studyframe}

\textbf{The chants} are Missal rubrics \rn{298}: all Sundays, of the I or the II class, have
their own Mass; \latin{Attamen dominicae post Epiphaniam, quae transferuntur inter dominicam
XXIII et XXIV post Pentecosten, sumunt antiphonas ad Introitum, ad Offertorium et ad
Communionem, necnon graduale et Alleluia cum suo versu a dominica XXIII post Pentecosten,
retentis orationibus, Epistola et Evangelio propriis} — however, the Sundays after the Epiphany
that are transferred between the Twenty-third and Twenty-fourth Sundays after Pentecost take
the Introit, Offertory, and Communion antiphons, and the Gradual and Alleluia with its verse,
from the Twenty-third Sunday after Pentecost, their own orations, Epistle, and Gospel being
retained.

\textbf{The trigger} is a rubric printed in the Missal itself at the end of the Twenty-third
Sunday after Pentecost:

\begin{studyframe}
\latin{Si dominicae post Pentecosten fuerint plures quam XXIV, tunc post XXIII resumuntur Missae
dominicarum, quae superfuerunt post Epiphaniam, ut infra habentur, iuxta ordinem qui in rubricis
invenitur. Et ultimo loco semper ponitur Missa dominicae XXIV, ut infra 422.}
\medskip

\textit{Working gloss.} If there are more than 24 Sundays after Pentecost, then after the
Twenty-third the Masses of the Sundays that were left over after the Epiphany are resumed, as
they are given below, according to the order found in the rubrics. And in the last place is
always put the Mass of the Twenty-fourth Sunday, as below at page 422.
\end{studyframe}

\subsection{What the three loci jointly establish}

\begin{enumerate}[leftmargin=1.6em, itemsep=0.35em]
\item \textbf{Resumption occurs only when $P > 24$.} The Missal's trigger says \latin{plures
quam XXIV}; \rn{18} enumerates only the cases 25, 26, 27, and 28. In a year of 24 Sundays after
Pentecost, nothing is resumed and the twenty-four printed formularies are used in order.
\item \textbf{The resumed Masses are taken from the top down.} In every enumerated case the
Sixth Sunday after the Epiphany is the last resumed, the Fifth the second-last, and so on. The
consequence is that the formularies dropped are always the \emph{lowest} unused ordinals, never
the highest.
\item \textbf{Only the Third through the Sixth are ever resumed.} Case \textit{d}, with 28
Sundays after Pentecost, reaches down to the Third; no enumerated case reaches the Second or the
First, so a Second Sunday after the Epiphany left unused is never resumed but simply omitted.
\item \textbf{The Twenty-fourth formulary is always last.} Both the code and the Missal say
\latin{semper}. It is not the twenty-fourth Sunday in the count that keeps this Mass; it is the
last Sunday, whatever its ordinal.
\item \textbf{The resumed Sunday keeps its own orations, Epistle, and Gospel, and borrows its
chants.} That is why the Missal prints four separate formularies — \latin{Dominica tertia},
\latin{quarta}, \latin{quinta}, and \latin{sexta quae superfuit post Epiphaniam}, each headed
\latin{II classis} — rather than instructing the celebrant to turn back to the Epiphany section.
Each printed resumed Mass already has the Twenty-third Sunday's Introit \latin{Dicit Dominus:
Ego cogito cogitationes pacis}, Gradual \latin{Liberasti nos, Domine}, Alleluia and Offertory
\latin{De profundis}, and Communion \latin{Amen dico vobis} joined to its own Epiphany orations
and readings.
\end{enumerate}

Point 3 needs one refinement that a purely arithmetical reading misses. In 2008, $E$ was 1: only
the First Sunday after the Epiphany occurred, so the Second through the Sixth were all unused,
five formularies in all. $P$ was 28, giving four resumption slots, which case \textit{d} fills
with the Third, Fourth, Fifth, and Sixth. The Second Sunday after the Epiphany therefore
occurred nowhere in 2008 — neither in January nor in November. It is dropped by the rule's own
top-down order, exactly as \rn{18}'s closing clause contemplates, and it is dropped precisely
because it is the lowest unused ordinal. Case \textit{d} does reach the Third; it is the
\emph{Second} that no case ever reaches.

\subsection{Why an arithmetical ordinal cannot identify the last Sunday}

The most consequential feature of this mechanism is the one hardest to see: after the
Twenty-third Sunday after Pentecost, the ordinal number of a Sunday and the ordinal number of
its Mass formulary come apart, and the last Sunday of the year is identified by \emph{position},
not by arithmetic.

Consider a year with 27 Sundays after Pentecost. Counting Trinity Sunday as the first, the
Sundays run to a twenty-seventh. But the Masses celebrated on the last four of them are, in
order: the Fourth after the Epiphany, the Fifth after the Epiphany, the Sixth after the
Epiphany, and the Twenty-fourth after Pentecost. The Sunday that is twenty-seventh in the count
uses a formulary printed as the twenty-fourth. Nothing in the year bears the title
``Twenty-seventh Sunday after Pentecost''; the Missal prints no such formulary, and \rn{18} does
not create one — it says that the twenty-seventh place \latin{erit quae inscribitur VI post
Epiphaniam}, will be the one inscribed as the Sixth after the Epiphany.

Hence the two propositions that a reference must keep apart:

\begin{studybox}{Ordinal and formulary}
\textbf{The ordinal} of a Sunday after Pentecost is its position in the count from Trinity
Sunday. It runs from 1 to $P$, and $P$ can be 23, 25, 26, 27, or 28 as well as 24.\\[0.35em]
\textbf{The formulary} used on that Sunday is the one \rn{18} assigns to that position. For the
first twenty-three positions it is the like-numbered Mass after Pentecost. For positions 24 and
beyond it is a resumed Epiphany Mass, except for the final position, which always takes the
Mass headed \latin{Dominica XXIV et ultima post Pentecosten}.
\end{studybox}

The Missal itself signals the distinction in its running head. The twenty-fourth formulary is
not titled \latin{Dominica XXIV post Pentecosten} but \latin{Dominica XXIV et ultima post
Pentecosten} — the Twenty-fourth \emph{and last}. The number is a place in the book; the word
\latin{ultima} is the operative identification. To say ``the last Sunday after Pentecost is the
twenty-fourth'' is true of the formulary and false of the count in about four years out of
five — the count is exactly twenty-four in only 44 of the 201 years from 1900 to 2100. To
compute the last Sunday's identity by adding one to twenty-three, or by counting weeks from
Pentecost, is to produce an ordinal for which the Missal has no Mass.

This is why an Ordo, and not a formula, is the competent statement of what is celebrated on a
given November Sunday, and why \rn{7}'s insistence that precedence be settled by the table and
not by inference has an arithmetical counterpart here: the identity of the last Sundays is
settled by the rubric's enumerated cases and not by counting.

\subsection{The year of twenty-three Sundays: a genuinely contested case}

Neither the 1960 code nor the 1962 Missal legislates for the case $P = 23$. \rn{18} enumerates
25, 26, 27, and 28; the Missal's trigger addresses \latin{plures quam XXIV}. A search of the
promulgated code, of the Missal's general rubrics and its Proper of the Season, and of the
pre-1955 general rubrics in the comparison witness consulted, found no provision using
\latin{pauciores}, ``fewer'', or any equivalent for the shortfall case. The gap is old: the
identified pre-1955 witness of Section~\ref{sec:witness} carries the resumption rubric in wording
identical to the 1962 Missal's apart from its internal page cross-reference (445 for 422) and the
grade given the resumed formularies (\latin{Semiduplex} for \latin{II classis}), and it too is
silent on the shortfall.

The case is rare but real. Between 1900 and 2100 it occurs four times: 1943, 2011, 2038, and
2095 — in each of which Easter falls on 24 or 25 April. The Missal's own \latin{Tabella
temporaria} prints 23 in the Sundays-after-Pentecost column for 2011, and its perpetual
\latin{Tabula paschalis antiqua reformata} prints 23 for both of the latest Easter rows. The
book therefore knows the case exists and does not say what to do in it.

Two readings are defensible on the text, and this reference gives both rather than choosing
silently.

\begin{studybox}{Reading A: the last place always takes the Twenty-fourth Mass}
\textbf{Claim.} In a year of twenty-three Sundays after Pentecost, the twenty-third and last
Sunday takes the Mass headed \latin{Dominica XXIV et ultima post Pentecosten}, and the
formulary of the Twenty-third Sunday is omitted that year.\\[0.4em]
\textbf{Textual support.} The clause \latin{Ultimo tamen loco semper ponitur ea quae in ordine
est XXIV post Pentecosten} is unrestricted in its own terms: \latin{semper} is not qualified by
the enumerated cases, and the subject is identified by its place in the printed order, not by
the ordinal of the Sunday. The formulary's own title, \latin{et ultima}, describes it as the
Mass of the last Sunday. The clause \latin{omissis, si opus sit, ceteris, quae aliquando locum
habere non possunt} — the rest being omitted, if need be, which sometimes cannot have a place —
reads naturally as covering whatever printed Mass finds no place, and in a twenty-three-Sunday
year the Mass that finds no place is the Twenty-third.\\[0.4em]
\textbf{Consequence.} The eschatological Gospel of the last Sunday, Matthew 24:15--35, is heard
every year without exception.
\end{studybox}

\begin{studybox}{Reading B: the count simply stops, and the Twenty-fourth Mass is omitted}
\textbf{Claim.} In a year of twenty-three Sundays after Pentecost, the Sundays are celebrated in
their printed order to the Twenty-third, which is the last; the Twenty-fourth formulary is not
used that year.\\[0.4em]
\textbf{Textual support.} Both loci embed the \latin{semper} clause inside a construction whose
premise is a surplus. \rn{18} governs \latin{dominicae post Epiphaniam quae \dots\ impediuntur}
and states the ``last place'' rule as a proviso — \latin{tamen} — on the four enumerated
resumption cases; the Missal's rubric states it inside a period beginning \latin{Si dominicae
post Pentecosten fuerint plures quam XXIV, tunc\dots}. On this reading \latin{ultimo loco}
means ``last among the resumed sequence'', and the rule is that resumed Epiphany Masses must not
be allowed to displace the Twenty-fourth from the end. Where no resumption occurs, the proviso
has nothing to govern. \rn{14}'s principle that an impeded Sunday's Mass is neither anticipated
nor resumed points the same way: a formulary for which the year has no Sunday is simply not
said.\\[0.4em]
\textbf{Consequence.} In roughly one year in fifty the twenty-fourth formulary is not used.
\end{studybox}

\textit{Editorial assessment, offered as such.} Reading A is the stronger of the two on the
wording, because \latin{semper} is an unusually emphatic word to place inside a subordinate
condition, because the formulary's printed title says \latin{et ultima}, and because the older
books gave the same Mass the same title under a rubrical tradition that plainly regarded it as
the Mass of the year's last Sunday. Reading B is not thereby refuted: its structural argument
about the scope of the proviso is sound Latin, and it has the advantage of requiring no
displacement of a printed formulary in a case the legislator never addressed.

\begin{studybox}{What would settle it}
A dated Ordo for 1943, 2011, 2038, or 2095, published under the 1960 rubrics by a competent
authority — a diocese, an institute using the older books, or a publisher whose Ordo carries an
ecclesiastical approbation — stating which Mass is assigned to the last Sunday after Pentecost
in that year. One such witness decides the practical question; a divergence between two would
establish that the question is genuinely open in practice as well as in text. This reference
locates and reads no such Ordo, and therefore records the case as unresolved rather than
resolving it by preference.
\end{studybox}

\subsection{The four resumed formularies as the Missal prints them}

\begin{center}
\small
\renewcommand{\arraystretch}{1.12}
\begin{tabular}{@{}>{\raggedright\arraybackslash}p{0.42\linewidth} >{\raggedright\arraybackslash}p{0.12\linewidth} >{\raggedright\arraybackslash}p{0.40\linewidth}@{}}
\toprule
\textbf{Printed heading} & \textbf{Class} & \textbf{Proper elements retained} \\
\midrule
\latin{Dominica tertia quae superfuit post Epiphaniam} & II & Orations, Epistle (Rom.\ 12:16--21), Gospel (Mt.\ 8:1--13) \\
\latin{Dominica quarta quae superfuit post Epiphaniam} & II & Orations, Epistle (Rom.\ 13:8--10), Gospel (Mt.\ 8:23--27) \\
\latin{Dominica quinta quae superfuit post Epiphaniam} & II & Orations, Epistle, Gospel \\
\latin{Dominica sexta quae superfuit post Epiphaniam} & II & Orations, Epistle, Gospel \\
\midrule
\multicolumn{3}{@{}>{\raggedright\arraybackslash}p{0.96\linewidth}@{}}{\textit{All four borrow the Introit, Gradual, Alleluia with its verse, Offertory, and Communion of the Twenty-third Sunday after Pentecost (Missal rubrics \rn{298}).}} \\
\bottomrule
\end{tabular}
\end{center}

That the Missal prints exactly four such formularies, and no First or Second, is itself the
book's confirmation that only the Third through the Sixth are ever resumed. The Third and Fourth
were read in the facsimile and their proper Epistles and Gospels are given above; the Fifth and
Sixth were located in the same run of pages but their readings were not separately collated for
this reference, and are therefore not cited here.
